\subsection{实验结论}
本实验利用 Python 程序对 DES 的三个 S 盒 $S_1,S_2,S_3$ 进行了差分分布分析，成功生成了对应的差分分布表，并统计了其主要安全性指标。实验结果表明：

\begin{enumerate}
    \item 三个 S 盒的最大差分传播概率均为 $0.25$；
    \item 三个 S 盒均具有 8 种非零差分传播概率；
    \item 不同 S 盒达到最大概率的差分对数量不同，说明其差分结构存在差异；
    \item 当输入差分非零时，输出差分仍可能为零，这是由 S 盒的非单射性质决定的。
\end{enumerate}
\par\indent 综合来看，DES 的 S 盒通过精心设计，使得差分分布尽量分散，避免出现过高的差分传播概率，从而增强了 DES 对差分密码分析的抵抗能力。本实验通过对 $S_1,S_2,S_3$ 的具体统计验证了这一点，也说明差分分布表是评估分组密码非线性部件安全性的重要工具。

\par\indent 通过本实验，进一步理解了差分密码分析的基本原理以及 DES S 盒在抗差分攻击设计中的重要作用。
